% MT028 / explanation — canonical editable content.
训练误差低表示模型能适应训练数据；测试误差高表示这套规律没有顺利推广。造成差距的原因可能在模型，也可能在数据来源。逐项看：



\begin{itemize}

\item \textbf{A选。} 训练集太小，学到的规律容易受到个别样本影响。例如只看到少数客户时，模型可能把某个碰巧共同的特征当成购买原因。

\item \textbf{B不选。} 测试集大只是让错误率估计通常更稳定，不会因为题数更多就使每题的预期错误率上升。若新增测试数据来自不同人群，原因应归到C。

\item \textbf{C选。} 训练与测试分布不同，即使模型把训练规律学得合理，换一个人群、设备或环境也可能失败。

\item \textbf{D选。} 模型表达能力过强，可能记住训练中的偶然细节，产生过拟合。可通过正则化或限制复杂度处理。

\item \textbf{E不选。} 模型太弱通常连训练数据也拟合不好，对应训练和测试都差。本题已经说明训练误差低；而且测试集用于检验推广能力，不是另一个要单独拟合的训练集。

\end{itemize}


处理A、D时可以补代表性数据或控制复杂度；处理C则要检查两批数据到底哪里不同。这个区别决定你下一步应该改数据还是改模型。


\textbf{出处：} 2021B*期中 Q2，原卷第2页（5分）


\textbf{答案依据：} 2021B*期中 Q2 随卷／配套参考答案；条件与解释经复核。
